% !atlas-meta
% id: yoga-sutras
% title: Yoga Sutras
% author: patanjali
% authorName: Patanjali
% language: English
% sourceLanguage: Sanskrit
% translator: Charles Johnston
% translationYear: 1912
% source: https://en.wikisource.org/wiki/Yoga_Sutras
% license: Public Domain (translator Charles Johnston d. 1931)
% status: complete
% !end-atlas-meta
\documentclass[12pt]{article}
\usepackage[utf8]{inputenc}
\usepackage{ebgaramond}
\usepackage[margin=1.2in]{geometry}
\usepackage{parskip}
\newcommand{\work}[1]{\begin{center}\LARGE\scshape #1\end{center}}
\newcommand{\attrib}[2]{\begin{center}\large #1\\[2pt]\small\itshape trans.\ #2\end{center}\vspace{1em}}
\newcommand{\para}[1][]{\par\noindent\ifx&#1&\else\textsuperscript{\footnotesize #1}~\fi}
\begin{document}
% !body-start
\work{The Yoga Sutras of Patanjali}
\attrib{Patanjali}{Charles Johnston (public domain)}

\section*{Book I}

\para[1] OM: Here follows Instruction in Union.

\para[2] Union, spiritual consciousness, is gained through control of the versatile psychic nature.

\para[3] Then the Seer comes to consciousness in his proper nature.

\para[4] Heretofore the Seer has been enmeshed in the activities of the psychic nature.

\para[5] The psychic activities are five; they are either subject or not subject to the five hindrances (Book II, 3).

\para[6] These activities are: Sound intellection, unsound intellection, predication, sleep, memory.

\para[7] The elements of sound intellection are: direct observation, inductive reason, and trustworthy testimony.

\para[8] Unsound intellection is false understanding, not resting on a perception of the true nature of things.

\para[9] Predication is carried on through words or thoughts not resting on an object perceived.

\para[10] Sleep is the psychic condition which rests on mind states, all material things being absent.

\para[11] Memory is holding to mind-images of things perceived, without modifying them.

\para[12] The control of these psychic activities comes through the right use of the will, and through ceasing from self- indulgence.

\para[13] The right use of the will is the steady, effort to stand in spiritual being.

\para[14] This becomes a firm resting-place, when followed long, persistently, with earnestness.

\para[15] Ceasing from self-indulgence is conscious mastery over the thirst for sensuous pleasure here or hereafter.

\para[16] The consummation of this is freedom from thirst for any mode of psychical activity, through the establishment of the spiritual man.

\para[17] Meditation with an object follows these stages: first, exterior examining, then interior judicial action, then joy, then realization of individual being.

\para[18] After the exercise of the will has stilled the psychic activities, meditation rests only on the fruit of former meditations.

\para[19] Subjective consciousness arising from a natural cause is possessed by those who have laid aside their bodies and been absorbed into subjective nature.

\para[20] For the others, there is spiritual consciousness, led up to by faith, valour right mindfulness, one-pointedness, perception.

\para[21] Spiritual consciousness is nearest to those of keen, intense will.

\para[22] The will may be weak, or of middle strength, or intense.

\para[23] Or spiritual consciousness may be gained by ardent service of the Master.

\para[24] The Master is the spiritual man, who is free from hindrances, bondage to works, and the fruition and seed of works.

\para[25] In the Master is the perfect seed of Omniscience.

\para[26] He is the Teacher of all who have gone before, since he is not limited by Time.

\para[27] His word is OM.

\para[28] Let there be soundless repetition of OM and meditation thereon.

\para[29] Thence come the awakening of interior consciousness, and the removal of barriers.

\para[30] The barriers to interior consciousness, which drive the psychic nature this way and that, are these: sickness, inertia, doubt, lightmindedness, laziness, intemperance, false notions, inability to reach a stage of meditation, or to hold it when reached.

\para[31] Grieving, despondency, bodily restless ness, the drawing in and sending forth of the life-breath also contribute to drive the psychic nature to and fro.

\para[32] Steady application to a principle is the way to put a stop to these.

\para[33] By sympathy with the happy, compassion for the sorrowful, delight in the holy, disregard of the unholy, the psychic nature moves to gracious peace.

\para[34] Or peace may be reached by the even sending forth and control of the life-breath.

\para[35] Faithful, persistent application to any object, if completely attained, will bind the mind to steadiness.

\para[36] As also will a joyful, radiant spirit.

\para[37] Or the purging of self-indulgence from the psychic nature.

\para[38] Or a pondering on the perceptions gained in dreams and dreamless sleep.

\para[39] Or meditative brooding on what is dearest to the heart.

\para[40] Thus he masters all, from the atom to the Infinite.

\para[41] When the perturbations of the psychic nature have all been stilled, then the consciousness, like a pure crystal, takes the colour of what it rests on, whether that be the perceiver, perceiving, or the thing perceived.

\para[42] When the consciousness, poised in perceiving, blends together the name, the object dwelt on and the idea, this is perception with exterior consideration.

\para[43] When the object dwells in the mind, clear of memory-pictures, uncoloured by the mind, as a pure luminous idea, this is perception without exterior or consideration.

\para[44] The same two steps, when referring to things of finer substance, are said to be with, or without, judicial action of the mind.

\para[45] Subtle substance rises in ascending degrees, to that pure nature which has no distinguishing mark.

\para[46] The above are the degrees of limited and conditioned spiritual consciousness, still containing the seed of separateness.

\para[47] When pure perception without judicial action of the mind is reached, there follows the gracious peace of the inner self.

\para[48] In that peace, perception is unfailingly true.

\para[49] The object of this perception is other than what is learned from the sacred books, or by sound inference, since this perception is particular.

\para[50] The impress on the consciousness springing from this perception supersedes all previous impressions.

\para[51] When this impression ceases, then, since all impressions have ceased, there arises pure spiritual consciousness, with no seed of separateness left.

\section*{Book II}

\para[1] The practices which make for union with the Soul are: fervent aspiration, spiritual reading, and complete obedience to the Master.

\para[2] Their aim is, to bring soul-vision, and to wear away hindrances.

\para[3] These are the hindrances: the darkness of unwisdom, self-assertion, lust hate, attachment.

\para[4] The darkness of unwisdom is the field of the others. These hindrances may be dormant, or worn thin, or suspended, or expanded.

\para[5] The darkness of ignorance is: holding that which is unenduring, impure, full of pain, not the Soul, to be eternal, pure, full of joy, the Soul.

\para[6] Self -assertion comes from thinking of the Seer and the instrument of vision as forming one self.

\para[7] Lust is the resting in the sense of enjoyment.

\para[8] Hate is the resting in the sense of pain.

\para[9] Attachment is the desire toward life, even in the wise, carried forward by its own energy.

\para[10] These hindrances, when they have become subtle, are to be removed by a countercurrent

\para[11] Their active turnings are to be removed by meditation.

\para[12] The burden of bondage to sorrow has its root in these hindrances. It will be felt in this life, or in a life not yet manifested.

\para[13] From this root there grow and ripen the fruits of birth, of the life-span, of all that is tasted in life.

\para[14] These bear fruits of rejoicing, or of affliction, as they are sprung from holy or unholy works.

\para[15] To him who possesses discernment, all personal life is misery, because it ever waxes and wanes, is ever afflicted with restlessness, makes ever new dynamic impresses in the mind; and because all its activities war with each other.

\para[16] This pain is to be warded off, before it has come.

\para[17] The cause of what is to be warded off, is the absorption of the Seer in things seen.

\para[18] Things seen have as their property manifestation, action, inertia. They form the basis of the elements and the sense-powers. They make for experience and for liberation.

\para[19] The grades or layers of the Three Potencies are the defined, the undefined, that with distinctive mark, that without distinctive mark.

\para[20] The Seer is pure vision. Though pure, he looks out through the vesture of the mind.

\para[21] The very essence of things seen is, that they exist for the Seer.

\para[22] Though fallen away from him who has reached the goal, things seen have not alto fallen away, since they still exist for others.

\para[23] The association of the Seer with things seen is the cause of the realizing of the nature of things seen, and also of the realizing of the nature of the Seer.

\para[24] The cause of this association is the darkness of unwisdom.

\para[25] The bringing of this association to an end, by bringing the darkness of unwisdom to an end, is the great liberation; this is the Seer's attainment of his own pure being.

\para[26] A discerning which is carried on without wavering is the means of liberation.

\para[27] His illuminations is sevenfold, rising In successive stages.

\para[28] From steadfastly following after the means of Yoga, until impurity is worn away, there comes the illumination of thought up to full discernment.

\para[29] The eight means of Yoga are: the Commandments, the Rules, right Poise, right Control of the life-force, Withdrawal, Attention, Meditation, Contemplation.

\para[30] The Commandments are these: nom injury, truthfulness, abstaining from stealing, from impurity, from covetousness.

\para[31] The Commandments, not limited to any race, place, time or occasion, universal, are the great obligation.

\para[32] The Rules are these: purity, serenity fervent aspiration, spiritual reading, and per feet obedience to the Master.

\para[33] When transgressions hinder, the weight of the imagination should be thrown' on the opposite side.

\para[34] Transgressions are injury, falsehood, theft, incontinence, envy; whether committed, or caused, or assented to, through greed, wrath, or infatuation; whether faint, or middling, or excessive; bearing endless, fruit of ignorance and pain. Therefore must the weight be cast on the other side.

\para[35] Where non-injury is perfected, all enmity ceases in the presence of him who possesses it.

\para[36] When he is perfected in truth, all acts and their fruits depend on him.

\para[37] Where cessation from theft is perfected, all treasures present themselves to him who possesses it.

\para[38] For him who is perfect in continence, the reward is valour and virility.

\para[39] Where there is firm conquest of covetousness, he who has conquered it awakes to the how and why of life.

\para[40] Through purity a withdrawal from one's own bodily life, a ceasing from infatuation with the bodily life of others.

\para[41] To the pure of heart come also a quiet spirit, one-pointed thought, the victory over sensuality, and fitness to behold the Soul.

\para[42] From acceptance, the disciple gains happiness supreme.

\para[43] The perfection of the powers of the bodily vesture comes through the wearing away of impurities, and through fervent aspiration.

\para[44] Through spiritual reading, the disciple gains communion with the divine Power on which his heart is set.

\para[45] Soul-vision is perfected through perfect obedience to the Master.

\para[46] Right poise must be firm and without strain.

\para[47] Right poise is to be gained by steady and temperate effort, and by setting the heart upon the everlasting.

\para[49] When this is gained, there follows the right guidance of the life-currents, the control of the incoming and outgoing breath.

\para[50] The life-current is either outward, or inward, or balanced; it is regulated according to place, time, number; it is prolonged and subtle.

\para[51] The fourth degree transcends external and internal objects.

\para[52] Thereby is worn away the veil which covers up the light.

\para[53] Thence comes the mind's power to hold itself in the light.

\para[54] The right Withdrawal is the disengaging of the powers from entanglement in outer things, as the psychic nature has been withdrawn and stilled.

\para[55] Thereupon follows perfect mastery over the powers.

\section*{Book III}

\para[1] The binding of the perceiving consciousness to a certain region is attention (dharana).

\para[2] A prolonged holding of the perceiving consciousness in that region is meditation (dhyana).

\para[3] When the perceiving consciousness in this meditative is wholly given to illuminating the essential meaning of the object contemplated, and is freed from the sense of separateness and personality, this is contemplation (samadhi).

\para[4] When these three, Attention, Meditation Contemplation, are exercised at once, this is perfectly concentrated Meditation (sanyama).

\para[5] By mastering this perfectly concentrated Meditation, there comes the illumination of perception.

\para[6] This power is distributed in ascending degrees.

\para[7] This threefold power, of Attention, Meditation, Contemplation, is more interior than the means of growth previously described.

\para[8] But this triad is still exterior to the soul vision which is unconditioned, free from the seed of mental analyses.

\para[9] One of the ascending degrees is the development of Control. First there is the overcoming of the mind-impress of excitation. Then comes the manifestation of the mind-impress of Control. Then the perceiving consciousness follows after the moment of Control.

\para[10] Through frequent repetition of this process, the mind becomes habituated to it, and there arises an equable flow of perceiving consciousness.

\para[11] The gradual conquest of the mind's tendency to flit from one object to another, and the power of one-pointedness, make the development of Contemplation.

\para[12] When, following this, the controlled manifold tendency and the aroused one-pointedness are equally balanced parts of the perceiving consciousness, his the development of one-pointedness.

\para[13] Through this, the inherent character, distinctive marks and conditions of being and powers, according to their development, are made clear.

\para[14] Every object has its characteristics which are already quiescent, those which are active, and those which are not yet definable.

\para[15] Difference in stage is the cause of difference in development.

\para[16] Through perfectly concentrated Meditation on the three stages of development comes a knowledge of past and future.

\para[17] The sound and the object and the thought called up by a word are confounded because they are all blurred together in the mind. By perfectly concentrated Meditation on the distinction between them, there comes an understanding of the sounds uttered by all beings.

\para[18] When the mind-impressions become visible, there comes an understanding of previous births.

\para[19] By perfectly concentrated Meditation on mind-images is gained the understanding of the thoughts of others.

\para[20] But since that on which the thought in the mind of another rests is not objective to the thought-reader's consciousness, he perceives the thought only, and not also that on which the thought rests.

\para[21] By perfectly concentrated Meditation on the form of the body, by arresting the body's perceptibility, and by inhibiting the eye's power of sight, there comes the power to make the body invisible.

\para[22] The works which fill out the life-span may be either immediately or gradually operative. By perfectly concentrated Meditation on these comes a knowledge of the time of the end, as also through signs.

\para[23] By perfectly concentrated Meditation on sympathy, compassion and kindness, is gained the power of interior union with others.

\para[24] By perfectly concentrated Meditation on power, even such power as that of the elephant may be gained.

\para[25] By bending upon them the awakened inner light, there comes a knowledge of things subtle, or concealed, or obscure.

\para[26] By perf ectly concentrated Meditation on the sun comes a knowledge of the worlds.

\para[27] By perfectly concentrated Meditation on the moon comes a knowledge of the lunar mansions.

\para[28] By perfectly concentrated Meditation on the fixed pole-star comes a knowledge of the motions of the stars.

\para[29] Perfectly concentrated Meditation on the centre of force in the lower trunk brings an understanding of the order of the bodily powers.

\para[30] By perfectly concentrated Meditation on the centre of f orce in the well of the throat, there comes the cessation of hunger and thirst.

\para[31] By perfectly concentrated Meditation on the centre of force in the channel called the "tortoise-formed," comes steadfastness.

\para[32] Through perfectly concentrated Meditation on the light in the head comes the vision of the Masters who have attained.

\para[33] Or through the divining power of tuition he knows all things.

\para[35] The personal self seeks to feast on life, through a failure to perceive the distinction between the personal self and the spiritual man. All personal experience really exists for the sake of another: namely, the spiritual man. By perfectly concentrated Meditation on experience for the sake of the Self, comes a knowledge of the spiritual man.

\para[36] Thereupon are born the divine power of intuition, and the hearing, the touch, the vision, the taste and the power of smell of the spiritual man.

\para[37] These powers stand in contradistinction to the highest spiritual vision. In mani- festation they are called magical powers.

\para[38] Through the weakening of the causes of bondage, and by learning the method of sassing, the consciousness is transf erred to the other body.

\para[39] Through mastery of the upward-life comes freedom from the dangers of water, morass, and thorny places, and the power of ascension is gained.

\para[40] By mastery of the binding-life comes radiance.

\para[41] From perfectly concentrated Meditation on the correlation of hearing and the ether, comes the power of spiritual hearing.

\para[42] By perfectly concentrated Meditation em the correlation of the body with the ether, and by thinking of it as light as thistle-down, will come the power to traverse the ether.

\para[43] When that condition of consciousness s reached, which is far-reaching and not con- fined to the body, which is outside the body and not conditioned by it, then the veil which conceals the light is worn away.

\para[44] Mastery of the elements comes from perfectly concentrated Meditation on their five forms: the gross, the elemental, the subtle, the inherent, the purposive.

\para[45] Thereupon will come the manifestation of the atomic and other powers, which are the endowment of the body, together with its unassailable force.

\para[46] Shapeliness, beauty, force, the temper of the diamond: these are the endowments of that body.

\para[47] Mastery over the powers of perception and action comes through perfectly concentrated Meditation on their fivefold forms; namely, their power to grasp their distinctive nature, the element of self-consciousness in them, their inherence, and their purposiveness.

\para[48] Thence comes the power swift as thought, independent of instruments, and the mastery over matter.

\para[49] When the spiritual man is perfectly disentangled from the psychic body, he attains to mastery over all things and to a knowledge of all.

\para[50] By absence of all self-indulgence at this point, when the seeds of bondage to sorrow are destroyed, pure spiritual being is attained.

\para[51] There should be complete overcoming of allurement or pride in the invitations of the different realms of life, lest attachment to things evil arise once more.

\para[52] From perfectly concentrated Meditatetion on the divisions of time and their succession comes that wisdom which is born of discernment.

\para[53] Hence comes discernment between things which are of like nature, not distinguished by difference of kind, character or position.

\para[54] The wisdom which is born of discerns ment is starlike; it discerns all things, and all conditions of things, it discerns without succession: simultaneously.

\para[55] When the vessture and the spiritual man are alike pure, then perfect spiritual life is attained.

\section*{Book IV}

\para[1] Psychic and spiritual powers may be inborn, or they may be gained by the use of drugs, or by incantations, or by fervour, or by Meditation.

\para[2] The transfer of powers from one venture to another comes through the flow of the natural creative forces.

\para[3] The apparent, immediate cause is not the true cause of the creative nature-powers; but, like the husbandman in his field, it takes obstacles away.

\para[4] Vestures of consciousness are built up in conformity with the Boston of the feel- ing of selfhood.

\para[5] In the different fields of manifestation, the Consciousness, though one, is the elective cause of many states of consciousness.

\para[6] Among states of consciousness, that which is born of Contemplation is free from the seed of future sorrow.

\para[7] The works of followers after Union make neither for bright pleasure nor for dark pain The works of others make for pleasure or pain, or a mingling of these.

\para[8] From the force inherent in works comes the manifestation of those dynamic mind images which are conformable to the ripening out of each of these works.

\para[9] Works separated by different nature, or place, or time, are brought together by the correspondence between memory and dynamic impression.

\para[10] The series of dynamic mind-images is beginningless, because Desire is everlasting.

\para[11] Since the dynamic mind-images are held together by impulses of desire, by the wish for personal reward, by the substratum of mental habit, by the support of outer things desired; therefore, when these cease, the self reproduction of dynamic mind-images ceases.

\para[12] The difference between that which is past and that which is not yet come, according to their natures, depends on the difference of phase of their properties.

\para[13] These properties, whether manifest or latent, are of the nature of the Three Potencies.

\para[14] The external manifestation of an object takes place when the transformations ore in the same phase.

\para[15] The paths of material things and of states of consciousness are distinct, as is manifest from the fact that the same object may produce different impressions in different minds.

\para[16] Nor do material objects defend upon a single mind, for how could they remain objective to others, if that mind ceased to think of them?

\para[17] An object is perceived, or not perceived, according as the mind is, or is not, tinged with the colour of the object.

\para[18] The movements of the psychic nature are perpetually ob jects of perception, since the Spiritual Man, who is the lord of them, remains unchanging.

\para[19] The Mind is not self-luminous, since it can be seen as an object.

\para[20] Nor could the Mind at the same time know itself and things external to it.

\para[21] If the Mind be thought of as seen by another more inward Mind, then there would be an endless series of perceiving Minds, and a confusion of memories.

\para[22] When the psychical nature takes on the form of the spiritual intelligence, by reflecting it, then the Self becomes conscious of its own spiritual intelligence.

\para[23] The psychic nature, taking on the colour of the Seer and of things seen, leads to the perception of all objects.

\para[24] The psychic nature, which has been printed with mind-images of innumerable material things, exists now f or the Spiritual Man, building for him.

\para[25] For him who discerns between the Mind and the Spiritual Man, there comes perfect fruition of the longing after the real being of the Self.

\para[26] Thereafter, the whole personal being bends toward illumination, toward Eternal Life.

\para[27] In the internals of the batik, other thoughts will arise, through the impressions of the dynamic mind-images.

\para[28] These are to be overcome as it was taught that hindrances should be overcome.

\para[29] He who, after he has attained, is wholly free from self, reaches the essence of all that can be known, gathered together like a cloud. This is the true spiritual consciousness.

\para[30] Thereon comes surcease from sorrow and the burden of toil.

\para[31] When all veils are rent, all stains washed away, his knowledge becomes infinite; little remains for him to know.

\para[32] Thereafter comes the completion of the series of transformations of the three nature potencies, since their purpose is attained.

\para[33] The series of transformations is divided into moments. When the series is completed, time gives place to duration.

\para[34] Pure spiritual life is, therefore, the in- verse resolution of the potencies of Nature, which have emptied themselves of their value for the Spiritual man; or it is the return of the power of pure Consciousness to its essential form.
% !body-end
\end{document}
